\section{Beurteilung des Kursverlaufs}
\label{sec:kurs-azn}

\subsection{Was die Reihe zeigt, und was an ihr bereinigt ist}

Bevor die Kurve gedeutet wird, geh\"ort eine Bereinigung offengelegt, die man
ihr nicht ansieht. Bis zum 2.~Februar 2026 wurde an der US-B\"orse kein Anteil
gehandelt, sondern ein Hinterlegungsschein \"uber eine halbe Stammaktie; seither
ist die Stammaktie selbst notiert.\quelleGB{2025}{223}\quelleMerken{s14gbxcacfxccd}
Die hier verwendete Reihe ist auf die Stammaktie zur\"uckgerechnet: der
Jahresschluss 2016 steht mit \USDje{\AznKursSchlussXVI}, w\"ahrend der
Hinterlegungsschein damals rund die H\"alfte davon kostete. W\"are die Reihe
nicht zur\"uckgerechnet, zeigte sie im Februar 2026 eine Verdopplung, die
keine war.

\"Uber die zehn Jahre bis zum \Datum{\AznKursTag} stieg der Kurs von
\USDje{\AznKursSchlussXVI} auf \USDje{\AznKurs}, das sind
\Pct{\AznRenditeZehn} pro Jahr, ohne Dividenden gerechnet und damit
untertrieben. Die H\"urde dieses Berichts, die Zehnjahresrendite des
Branchenindex, liegt bei \Pct{\Huerde}. Der Titel hat den Index in diesem
Zeitraum also geschlagen, und zwar bevor die Dividende hinzukommt.

\subsection{Was zu erkl\"aren ist: der R\"uckgang seit Februar 2026}

Zu erkl\"aren ist nicht der Aufstieg. Am \Datum{\AznKurstagHoch} stand die Aktie bei
\USDje{\AznKursHochXXVI}; heute steht sie \PctBetrag{\AznAbstandHoch} darunter
und \PctBetrag{\AznRenditeLaufend} unter dem Jahresschluss 2025. Das ist der
Vorgang, der eine Erkl\"arung verlangt, und zwar in einem Jahr, dessen erstes Halbjahr
\MioUSD{\AznUmsatzHJ} Umsatz brachte und in dem das Unternehmen seine Prognose
unver\"andert best\"atigt hat.\quelleErneut{s14hjxb}

\bk{Drei Gr\"o{\ss}en aus diesem Bericht kommen als Ursache in Betracht, und
sie lassen sich der Reihe nach pr\"ufen.}

\subsection{Der Fr\"uhindikator, nicht die Gewinn- und Verlustrechnung}

\bk{Die Ergebnisrechnung erkl\"art den R\"uckgang nicht. Sie zeigt das
Gegenteil: den h\"ochsten Umsatz, das h\"ochste Betriebsergebnis und den
h\"ochsten freien Mittelzufluss der Reihe (Tabelle~\ref{tab:azn-reihe}).
Was f\"allt, f\"allt gegen die ver\"offentlichten Zahlen, nicht mit ihnen.}

\bk{Der Fr\"uhindikator erkl\"art ihn besser. Farxiga, das gr\"o{\ss}te
Einzelmedikament mit \Pct{\AznFarxigaAnteil} des Konzernumsatzes, hat im
zweiten Quartal 2026 in den Vereinigten Staaten den Patentschutz verloren; der
Halbjahresumsatz des Produkts liegt hochgerechnet \Pct{\AznFarxigaHJWachstum}
unter dem Vorjahr, in den Vereinigten Staaten deutlicher.\quelleHJ{9}\quelleMerken{s15hjxj}
Der Verlust war seit dem Bericht 2025 angek\"undigt und ist nun eingetreten.
Ein Markt, der einen angek\"undigten Patentablauf beim Eintreten noch einmal
einpreist, verhaelt sich nicht widerspruechlich: Angek\"undigt war der
Zeitpunkt, nicht die Tiefe.}

\bk{Die dritte Gr\"o{\ss}e ist die Bewertung selbst, und sie ist von den
ersten beiden zu trennen. Ein Kurs, der f\"allt, w\"ahrend Umsatz und Ergebnis
steigen, ist ein Bewertungsereignis: Das Kurs-Gewinn-Verh\"altnis auf die
bereinigte Gr\"o{\ss}e ging von \Faktor{\AznKgvHoch} am Jahreshoch auf
\Faktor{\AznKursGewinnCore} zur\"uck, ohne dass sich der Nenner verschlechtert
h\"atte. Das ist wichtig f\"ur Teil~\ref{sec:flow-azn}: Ein Bewertungsr\"uckgang
geh\"ort nicht in die Ertragsszenarien, sonst wird er zweimal gez\"ahlt.}

\subsection{Was der Kurs nicht abbildet}

\bk{Zwei Dinge bleiben offen, und beide sind aus diesem Bericht nicht zu
beantworten. Erstens ist unbekannt, wie viel Umsatz insgesamt an auslaufenden
Patenten h\"angt; das Unternehmen ver\"offentlicht die Aufstellung
gesondert, und sie liegt hier nicht vor (Abschnitt~\ref{sec:loe-azn}).
Zweitens l\"asst sich der Anteil des R\"uckgangs, der auf den Sektor und nicht
auf das Unternehmen entf\"allt, mit den Mitteln dieses Berichts nicht sauber
abspalten; die H\"urde ist eine Zehnjahresgr\"o{\ss}e und kein Tagesindex.}
